\section{Content-verified retrieval}
\label{sec:retrieval}
\label{sec:lanes}

Content governs admission lives between candidate discovery and answering:
the Search, Librarian and Hybrid lanes use different retrieval affordances but
share one server-side content gate. The gate fetches document text, scores its
segments on a 0--8 rubric (topic 0--3, scenario 0--3, evidence 0--2), and
fails closed on missing text, tool errors or invalid verdicts.

\subsection{Gate contract and retrieval lanes}
\label{sec:trace}
\label{sec:gate}
\label{sec:gate-primitive}
\label{sec:twostage}

The chat endpoint implements two policies. Lane A uses a wide dense search
(top 30, cosine threshold 0.35) within a pinned library and may escalate once
when no verified evidence survives. Lane B navigates library and document
descriptions, then uses a deterministic scanner for literal or enumeration
questions and the judge for semantic questions; a fixed call budget bounds the
search, and exhaustion produces a scope and blind-spot report. Lane C is a
composed configuration used in evaluation, not a third built-in server policy:
it appends a hybrid directive to the chat prompt that can call vector and
catalog tools, merge up to twelve candidates, judge them in batches of six, and
return up to five survivors with bounded evidence. Table~\ref{tab:lanes} and
Fig.~\ref{fig:lanes} summarize the policies and shared answer contract.

The gate is the local Laya evidence classifier \cite{laya}, exposed as
\texttt{kb\_laya\_judge}. In reference mode, the server fetches each document,
judges its segments and returns provenance with a bounded evidence pack. The
verdict controls admission, not truth; the verifier remains subject to
calibration and transport limitations.

\begin{table*}[t]
\centering
\small
\setlength{\tabcolsep}{5pt}
\caption{Retrieval lanes as operating policies over a shared candidate and verification tool layer. Budgets describe the configured interface; latency and effectiveness are evaluated separately in Section~\ref{sec:eval}.}
\label{tab:lanes}
\begin{tabular}{@{}llll@{}}
\toprule
Lane & Retrieval affordance & Budget & Terminal behavior \\
\midrule
Search & Pinned-library dense search & 20 turns & Escalate once or answer/not-found \\
Librarian & Catalog navigation and query-shaped scan/judging & 12 calls & Report scope and blind spots \\
Hybrid & Server-side vector + catalog composition & 1 retrieval call & Shared verification, then answer/not-found \\
\bottomrule
\end{tabular}
\end{table*}

The two-stage tool builds a candidate pool with lexical search and optional
graph expansion, then applies dense ranking. It does not call the
content gate; a ranked chunk is not admissible evidence until verified, and a
similarity threshold is not an abstention mechanism. The 13 October catalog
incident showed that oversized tool returns can conceal documents before
reasoning begins. The tested repair adds a server-side character budget,
grouped views and pagination; its behavior is reported in the Evaluation
section.

\begin{figure*}[t]
\centering
\includegraphics[width=\textwidth]{fig2_lanes}
\caption{The three lanes reach one content-verification gate by different
routes: a pinned wide dense net that may escalate once when nothing survives,
a catalogue ladder whose query shape selects a literal scanner or the judge,
and a single server-side vector-and-catalogue call that merges and caps its
candidates at twelve. The gate reads the fetched text and fails closed;
zero verified evidence or an exhausted budget produces an explicit
not-found report naming the scope searched. Lane budgets are given in
Table~\ref{tab:lanes}.}
\label{fig:lanes}
\end{figure*}

\subsection{Honest abstention}
A candidate list may contain high-similarity but unsupported documents. The
answer contract therefore requires a not-found response when no segment meets
the evidence criterion, naming the searched libraries and known blind spots
rather than falling back to parametric knowledge; results and limits appear in
Section~\ref{sec:eval}.
